\documentclass[11pt,a4paper]{article}
\usepackage[utf8]{inputenc}
\usepackage[T1]{fontenc}
\usepackage[a4paper,margin=1in]{geometry}
\usepackage{longtable}
\usepackage{array}
\usepackage{hyperref}
\usepackage{xcolor}
\usepackage{enumitem}
\usepackage{titlesec}
\usepackage{fancyhdr}
\usepackage{parskip}
\usepackage{tabularx}
\usepackage{booktabs}

\hypersetup{
  colorlinks=true,
  linkcolor=black,
  urlcolor=blue,
  pdftitle={Morfeas WEB Help Manual},
  pdfauthor={Morfeas Project}
}

\pagestyle{fancy}
\fancyhf{}
\lhead{Morfeas WEB}
\chead{Help Manual}
\rhead{\thepage}
\setlength{\headheight}{14pt}
\setlist[itemize]{leftmargin=1.4em}
\setlist[enumerate]{leftmargin=1.6em}

\titleformat{\section}{\large\bfseries}{}{0pt}{}
\titleformat{\subsection}{\normalsize\bfseries}{}{0pt}{}
\titleformat{\subsubsection}{\normalsize\bfseries}{}{0pt}{}

\newcommand{\screenshotplaceholder}[2]{%
\begin{center}%
\fbox{\parbox{0.9\textwidth}{\textbf{Screenshot Placeholder: #1}\\#2}}%
\end{center}%
}

\newcommand{\messageexample}[2]{%
\noindent\textbf{Example Message (#1):}\\%
\begin{quote}\raggedright #2\end{quote}%
}

\title{Morfeas WEB\\Help Manual}
\author{}
\date{Current New Web UI}

\begin{document}
\maketitle
\thispagestyle{empty}
\tableofcontents
\newpage

\section{Introduction}
Morfeas WEB is the browser user interface for daily operation and configuration of the Morfeas system.
This manual describes the current \textbf{new} Morfeas web interface and is intended for operators and service users.

Unless a site-specific deployment says otherwise, open the system in a browser with:

\begin{itemize}
\item \texttt{http://<device-ip>/}
\end{itemize}

This manual does not describe the retired legacy web layout.

\screenshotplaceholder{Main Landing Page}{Insert a current screenshot of the ISO Channels Linker landing page, showing the top bar, menu button, search bar, and the main channel table.}

\section{Main Page}
The landing page is \textbf{ISO Channels Linker}. It is the main operational screen.

\subsection{What You See}
The page contains:

\begin{itemize}
\item a top bar with the FTP backup directory chip and rolling system summary,
\item a search bar,
\item action buttons for importing and adding channels,
\item the main ISO channel table,
\item a right-click context menu for channel actions,
\item a menu button that opens all secondary pages.
\end{itemize}

\subsection{Main Table}
The table shows the configured ISO channels and related runtime information.
Typical columns include:

\begin{itemize}
\item status,
\item ISO channel name,
\item device type,
\item connection anchor,
\item current value,
\item history,
\item next calibration,
\item minimum and maximum values,
\item description.
\end{itemize}

The page refreshes automatically while it is visible, so operators normally do not need to reload manually.

\subsection{Search}
Use \textbf{Search Input} to filter channels in the current table.

\subsection{Context Menu}
Right-click a row to access actions such as:

\begin{itemize}
\item Edit,
\item Delete,
\item Scale,
\item Calibration,
\item Export,
\item device-specific actions such as Replace or Replace TC16 when the selected channel supports them.
\end{itemize}

\section{Menu Overview}
The main menu is grouped as follows.

\begin{longtable}{|p{0.28\textwidth}|p{0.62\textwidth}|}
\hline
\textbf{Menu Item} & \textbf{Purpose} \\
\hline
Network Configuration & Edit hostname and Ethernet IPv4 settings. \\
\hline
Add Devices & View current devices and add or remove supported devices. \\
\hline
Backup \& Restore Channels & Configure FTP backup target, create backups, and restore saved bundles. \\
\hline
Check System Status & View live runtime details, log files, and system journal output. \\
\hline
Advanced Settings & View MAC/CAN information, configure NTP server, and manage ISOstandard XML. \\
\hline
System Update & Check for available Morfeas web updates and apply them when desired. \\
\hline
System Reboot & Reboot the device. \\
\hline
System Shutdown & Power off the device. \\
\hline
Version & Show current web and core version information. \\
\hline
Help & Download this manual as PDF. \\
\hline
\end{longtable}

\section{Network Configuration}
Open \textbf{Network Configuration} from the menu to edit hostname and Ethernet settings.

\subsection{Fields}
The page contains:

\begin{itemize}
\item \textbf{Hostname}
\item \textbf{Mode}: Static or DHCP
\item \textbf{Network IP}
\item \textbf{Gateway}
\item \textbf{DNS addr}
\end{itemize}

Preset buttons are available at the top of the page for quick loading of predefined values.

\subsection{Buttons}
\begin{itemize}
\item \textbf{Load Last}: reload the last confirmed configuration into the form.
\item \textbf{Set}: apply the current values.
\item \textbf{Close}: close the popup window.
\end{itemize}

\subsection{Operator Notes}
\begin{itemize}
\item If you change the device IP address, reconnect the browser to the new HTTP address after the change is applied.
\item DHCP mode disables manual IPv4 entry.
\item Presets are convenience values only. Review them before applying.
\end{itemize}

\screenshotplaceholder{Network Configuration Page}{Insert a screenshot showing Hostname, Mode, Network IP, Gateway, DNS addr, and the Load Last / Set / Close buttons.}

\messageexample{Progress}{Applying configuration: validating and writing system network settings...}
\messageexample{Success}{Configuration applied. If the IP changed, reconnect to the new device address in your browser.}
\messageexample{Failure}{Apply failed: <error details>}

\section{Add Devices}
Open \textbf{Add Devices} to manage the current device inventory.

\subsection{Devices Table}
The upper table shows the current devices with columns such as:

\begin{itemize}
\item Bus,
\item Device Type,
\item Device Name or Address,
\item IP Address.
\end{itemize}

\subsection{Toolbar}
\begin{itemize}
\item \textbf{Refresh}: reload the current device list.
\item \textbf{Remove}: remove selected removable devices.
\item \textbf{Add...}: open the add form.
\end{itemize}

\subsection{Add Form}
The add form lets the user enter:

\begin{itemize}
\item device type,
\item CAN interface,
\item device name,
\item additional device properties required by that device family.
\end{itemize}

\subsection{Operator Notes}
\begin{itemize}
\item Some devices, such as automatically detected SDAQ devices, may not be removable from this page.
\item The popup refreshes while active, and pauses when the popup is no longer in focus.
\item If the remove action is not available or shows a warning, review whether the selected device is auto-detected and protected from removal.
\end{itemize}

\section{Channel Replace Workflows}
Morfeas WEB supports both single-channel replacement and grouped TC16 replacement.

\subsection{Single-Channel Replace}
Single-channel replace uses the standard Edit Link popup in replace mode.

\subsubsection{Typical Flow}
\begin{enumerate}
\item In the main ISO channel table, right-click the channel you want to replace.
\item Choose the standard replace action.
\item The popup opens with the source type locked and the target path chosen through device search.
\item Select the replacement device path.
\item Review ISO code, description, min/max, alarms, unit, and calibration fields.
\item Click \textbf{Save}.
\end{enumerate}

\subsubsection{Type Safety}
The replace flow enforces compatibility rules.

\screenshotplaceholder{Single-Channel Replace Popup}{Insert a screenshot of the replace popup showing the locked source type, search button, ISO code, limits, alarms, and Save button.}

\begin{itemize}
\item Family mismatch is blocked.
\item If the source SDAQ subtype is known, subtype mismatch is blocked.
\item If the source SDAQ subtype is unknown, replacement is still allowed, but the user is warned to replace within the same type.
\end{itemize}

\subsubsection{Offline Type Display}
When an SDAQ channel has been identified before and later goes offline, the UI can display the last known type in a stale form such as:

\begin{itemize}
\item \texttt{SDAQ-TC1 (last known)}
\end{itemize}

If the subtype has never been identified, the type may be shown as unknown.

\messageexample{Warning}{Subtype unknown, make sure replace within same type.}
\messageexample{Blocked Replace}{Subtype mismatch: expected SDAQ-TC1, got SDAQ-AI8.}

\subsection{TC16 Group Replace}
The dedicated \textbf{Replace TC16} flow is separate from single-channel replace.
It is intended for a full 16-channel SDAQ-TC16 source group.

\subsubsection{When the Option Appears}
Replace TC16 only appears when the selected source channel belongs to a resolvable offline TC16 group that is complete as channels CH1 through CH16.
If the group is partial, the option is not shown and normal single-channel replace must be used.

\subsubsection{How It Works}
\begin{enumerate}
\item Right-click any channel belonging to the full TC16 source group.
\item Choose \textbf{Replace TC16}.
\item The UI highlights and selects all 16 source channels.
\item A dedicated Replace TC16 popup opens.
\item The popup shows the source group summary and the source channel mapping.
\item Choose one eligible unlinked TC16 target device.
\item Confirm the operation.
\item All 16 source channels are remapped to the target device in one atomic operation.
\end{enumerate}

\subsubsection{Important Notes}
\begin{itemize}
\item Replace TC16 is all-or-nothing.
\item If any validation fails, no channel mapping is changed.
\item Target selection is at device-group level, not one channel at a time.
\item If no eligible target is listed, there is currently no fully free TC16 target device available.
\end{itemize}

\screenshotplaceholder{Replace TC16 Popup}{Insert a screenshot showing the source group summary, CH1..CH16 mapping chips, available target list, and the Apply Replace TC16 button.}

\messageexample{Success}{Replace TC16 complete: 16 channels updated.}
\messageexample{Failure}{Failed to apply TC16 replace. [tc16\_target\_not\_unlinked]}

\section{Backup and Restore Channels}
Open \textbf{Backup \& Restore Channels} for FTP-based backup and restore.

\subsection{FTP Server Setup}
The page contains:

\begin{itemize}
\item \textbf{Host IP}
\item \textbf{Engine Number}
\item \textbf{Connect \& Fetch}
\item \textbf{Disconnect}
\end{itemize}

The default host value is prefilled for the expected FTP server. Credentials are stored server-side and are not edited in the UI.

\subsection{Backup to FTP}
Use \textbf{Backup Now} to create and upload a backup bundle for the current configuration.

\subsection{Restore from FTP}
Select a backup file from the list and click \textbf{Restore Selected}.

\subsection{Step-by-Step Example}
\subsubsection{Connect and Load Backup List}
\begin{enumerate}
\item Open \textbf{Backup \& Restore Channels}.
\item Confirm or enter the FTP server IP.
\item Enter the engine number.
\item Click \textbf{Connect \& Fetch}.
\item Wait for the connection result.
\item If the connection succeeds, the backup list is loaded for that engine directory.
\end{enumerate}

\subsubsection{Create a Backup}
\begin{enumerate}
\item Make sure the FTP connection is active.
\item Click \textbf{Backup Now}.
\item Wait for the success message.
\item Confirm that the new file appears in the backup list.
\end{enumerate}

\subsubsection{Restore a Backup}
\begin{enumerate}
\item Make sure the FTP connection is active.
\item Select one file from the backup list.
\item Click \textbf{Restore Selected}.
\item Wait for the restore result message.
\end{enumerate}

\subsection{Operator Notes}
\begin{itemize}
\item Backup and restore both require a reachable FTP server.
\item If the server is unavailable, the page shows an error and no system state is changed.
\item The FTP directory chip on the main page shows the currently connected engine directory when configuration exists.
\item If two sessions edit the FTP configuration, the page can detect configuration updates from the other session.
\end{itemize}

\screenshotplaceholder{FTP Backup and Restore Page}{Insert a screenshot showing the FTP Server Setup, Backup Now, Restore Selected, and the available backup list.}

\messageexample{Connect Success}{Connected to FTP at 10.193.135.70. Configuration saved.}
\messageexample{Connect Failure}{Connection failed: FTP connect failed}
\messageexample{Backup Success}{Backup complete}
\messageexample{Restore Failure}{Restore failed: <error details>}

\section{Check System Status}
Open \textbf{Check System Status} to inspect runtime data and logs.

\subsection{Tabs}
The page has three tabs.

\subsubsection{Details}
This tab shows live system and bus information in tables. Depending on the device and services, values may include:

\begin{itemize}
\item CPU temperature,
\item uptime,
\item bus utilization,
\item bus error rate,
\item detected SDAQ count,
\item incomplete SDAQ count,
\item bus voltage,
\item bus amperage,
\item shunt temperature,
\item last calibration timestamps.
\end{itemize}

\subsubsection{System Logs}
This tab is used for Morfeas application logger files.

\begin{itemize}
\item \textbf{Preview Logger}: choose one log file to preview.
\item \textbf{Reload}: reload the logger list.
\item \textbf{Export Logs}: select one or more log files and export them as one combined text file.
\end{itemize}

The preview area shows the selected file content in a terminal-style viewer.

\subsubsection{System Journal}
This tab is used for journal output from system services.

\begin{itemize}
\item \textbf{Scope}: choose all services, Morfeas system, Apache2, or SSH.
\item \textbf{Lines}: choose how many lines to load.
\item \textbf{Reload}: reload the journal view.
\item \textbf{Export Journal}: export the currently requested journal text.
\end{itemize}

\screenshotplaceholder{System Status Tabs}{Insert a screenshot showing the Details, System logs, and System Journal tabs with one example content area visible.}

\subsection{How to Use the Tabs}
\subsubsection{Details Tab}
Use this tab first when checking current runtime health.
If a service is active but values look stale, compare with the Logs or Journal tabs.

\subsubsection{System Logs Tab}
Use this tab for Morfeas component logs when you want application-level detail such as handler messages, update logs, backup logs, or runtime component reports.

\subsubsection{System Journal Tab}
Use this tab when the problem looks system-level rather than application-level, for example:

\begin{itemize}
\item service startup failures,
\item Apache problems,
\item SSH availability,
\item reboot or shutdown investigation.
\end{itemize}

\section{Advanced Settings}
Open \textbf{Advanced Settings} for machine information, NTP, and ISOstandard management.

\subsection{Machine Identifier}
Displays the MAC address.

\subsection{SDAQNet (CAN-IFs)}
Displays the current CAN interface information for \texttt{can0} and \texttt{can1}. These values are informational and read-only on this page.

\subsection{Time Synchronization}
The page shows the current NTP server and allows entering a new IPv4 NTP server. Use \textbf{Apply} to store the new value.

\messageexample{Apply Success}{Applied NTP server: 192.168.1.10}
\messageexample{Validation Error}{NTP server must be a valid IPv4 address.}

\subsection{ISOstandard}
This section lets the user:

\begin{itemize}
\item view the currently selected ISOstandard source,
\item select a different available ISOstandard file,
\item upload a new ISOstandard XML file,
\item review the loaded ISO point definitions in a table.
\end{itemize}

\section{System Update}
Open \textbf{System Update} to manage Morfeas web software updates.

\subsection{Current Behavior}
The page supports:

\begin{itemize}
\item checking whether an update is available,
\item applying the update,
\item deferring the update,
\item reloading the page after service recovery.
\end{itemize}

\subsection{Reminder Dot}
A red dot on the menu item indicates that the system has recorded an available update. The reminder stays visible until a later check confirms that no update is pending, or until a successful update clears it.

\subsection{Important Note}
The current update flow is intended for the Morfeas web deployment path. Core update is handled separately.

\screenshotplaceholder{System Update Page}{Insert a screenshot showing Check Again, Update Now, Update Later, Reload Page, and the status box.}

\messageexample{Update Available}{Update available. Click "Update Now" to apply, or "Update Later" to postpone.}
\messageexample{Update Success}{Update flow completed and web service is reachable. Refresh page (Ctrl+F5) to load latest files.}
\messageexample{Update Failed}{System update failed: <error details>}

\section{System Reboot and System Shutdown}
The system control pages provide single-confirmation actions for reboot and shutdown.

\subsection{System Reboot}
Use \textbf{Reboot Now} when you want to restart the device and return to service after boot. The page shows post-action guidance so the operator knows when to reconnect and refresh the browser.

\messageexample{Reboot Accepted}{Reboot command accepted. Wait about 60--90 seconds, then open the device URL again and refresh the browser.}

\subsection{System Shutdown}
Use \textbf{Shutdown Now} when you want to power off the device. After shutdown, the device must be powered on again manually before the web interface becomes available.

\messageexample{Shutdown Accepted}{Shutdown command accepted. Device will go offline. Power it on manually when needed.}

\section{Version}
The \textbf{Version} page shows version and commit information for:

\begin{itemize}
\item the Morfeas web repository,
\item the Morfeas core repository,
\item the related commit dates.
\end{itemize}

Use this page when confirming what software is installed on a device.

\section{Troubleshooting}
This section gives quick operator guidance for common problems.

\screenshotplaceholder{Troubleshooting Section Reference}{Optionally insert a screenshot collage of the key popups or status messages that operators most commonly use during troubleshooting.}

\subsection{The Web Page Does Not Open}
\begin{itemize}
\item Confirm that you are using the current device IP address over HTTP.
\item If the IP was recently changed, reconnect to the new address.
\item If the device was rebooted or updated, wait for services to come back and refresh the browser.
\item If still unavailable, check \textbf{System Journal} or use SSH to verify Apache and Morfeas services.
\end{itemize}

\subsection{FTP Backup Cannot Connect}
\begin{itemize}
\item Check that the FTP server IP is reachable from the device.
\item Confirm that the engine number is correct.
\item If the page says configuration was removed, reconnect first.
\item If the internal FTP server is currently unavailable, backup and restore cannot be tested end-to-end.
\end{itemize}

\subsection{No Backup Files Are Listed}
\begin{itemize}
\item Verify that the selected engine directory contains backup files.
\item Confirm that the FTP configuration belongs to the expected engine number.
\item Reload after reconnecting to the FTP server.
\end{itemize}

\subsection{Replace or Replace TC16 Option Is Missing}
\begin{itemize}
\item Confirm that the selected row supports that replace mode.
\item For Replace TC16, confirm that the source is a complete offline TC16 group with channels CH1 through CH16.
\item If the source group is partial, use normal single-channel replace instead.
\end{itemize}

\subsection{No Eligible TC16 Target Is Listed}
\begin{itemize}
\item No full unlinked TC16 target device is currently available.
\item If channels on the target are already linked, the target is blocked.
\item Refresh the popup after any device or link changes.
\end{itemize}

\subsection{Logs or Journal Look Empty}
\begin{itemize}
\item Use \textbf{Reload} once to refresh the selected tab.
\item Verify that the correct logger file or journal scope is selected.
\item If the service has not produced messages yet, the view may be temporarily empty.
\end{itemize}

\subsection{Update Reminder Dot Is Still Visible}
\begin{itemize}
\item The reminder is based on the stored update flag.
\item Run \textbf{Check Again} in the System Update page to refresh the latest status.
\item After a successful update or a clean re-check, the reminder should clear.
\end{itemize}

\section{Help PDF Maintenance}
The Help menu downloads a PDF file from:

\begin{itemize}
\item \texttt{docs/help\_manual.pdf}
\end{itemize}

The maintained source for this file is located under:

\begin{itemize}
\item \texttt{Morfeas\_WEB/docs/manual/}
\end{itemize}

This keeps the active user manual next to the current web code instead of inside the archived legacy documentation tree.

\section{Legacy Documentation Status}
The old documentation tree under \texttt{LOG-web/Docs/Morfeas\_WEB\_Docs} is retained only as historical reference. Its structure and API descriptions belong to the legacy implementation and should not be used as the active operator manual for the current web UI.

\end{document}
